% GENERATED FILE. Edit canonical lessons, then run npm run generate:books.
% canonical-source-hash: fnv1a64:c61586042cb6ac3d
% canonical-lessons: BN-C161-bhodro, BN-C161-komol, BN-C161-kora, BN-C161-sukhi, BN-C161-birokto

\chapter{Polite, Gentle, Strict, Contented, Annoyed}
\label{ch:bn-polite-gentle-strict-contented-annoyed}
\hlchaptermodality{161}

\begin{chapteropening}
\textbf{By the end of this chapter:} \emph{I can say polite, gentle, strict, contented and annoyed.}

Complete the last lesson of chapter 161: I can say polite, gentle, strict, contented and annoyed.
\end{chapteropening}

\section[\texorpdfstring{\bn{ভদ্র} (bhôdrô)}{ভদ্র (bhôdrô)}]{\bn{ভদ্র} (bhôdrô) — polite}

\label{lesson:BN-C161-bhodro}



\begin{quote}
Before the new one: say the Bengali for worried, then the Bengali for elderly.
\end{quote}

\subsection*{You'll want to know: \bn{ভদ্র}}
\textbf{\bn{ভদ্র}} — \emph{bhôdrô} — \textquotedblleft{}polite\textquotedblright{}.

\begin{grammarlens}[title={What you've built}]
One more word for this chapter.
\end{grammarlens}

\subsection*{Your turn}
\begin{itemize}
\raggedright
  \item \emph{Say it:} \emph{bhôdrô}
  \item \emph{Say it:} \emph{bhôdrô}, once more
  \item \emph{Say it:} say \emph{bhôdrô}
  \item \emph{Recall:} say the Bengali for worried, then the Bengali for elderly, then say \emph{bhôdrô} again
\end{itemize}

\begin{tcolorbox}[breakable,colback=teal!4,colframe=teal!35!black,title={Before you move on}]
What does \bn{ভদ্র} mean? (\textquotedblleft{}Polite\textquotedblright{}.) Say it once more.
\end{tcolorbox}
\section[\texorpdfstring{\bn{কোমল} (komôl)}{কোমল (komôl)}]{\bn{কোমল} (komôl) — gentle, tender}

\label{lesson:BN-C161-komol}



\begin{quote}
Before the new one: say the Bengali for elderly, then the Bengali for polite.
\end{quote}

\subsection*{You'll want to know: \bn{কোমল}}
\textbf{\bn{কোমল}} — \emph{komôl} — \textquotedblleft{}gentle, tender\textquotedblright{}.

\begin{grammarlens}[title={What you've built}]
One more word for this chapter.
\end{grammarlens}

\subsection*{Your turn}
\begin{itemize}
\raggedright
  \item \emph{Say it:} \emph{komôl}
  \item \emph{Say it:} \emph{komôl}, once more
  \item \emph{Say it:} say \emph{komôl}
  \item \emph{Recall:} say the Bengali for elderly, then the Bengali for polite, then say \emph{komôl} again
\end{itemize}

\begin{tcolorbox}[breakable,colback=teal!4,colframe=teal!35!black,title={Before you move on}]
What does \bn{কোমল} mean? (\textquotedblleft{}Gentle, tender\textquotedblright{}.) Say it once more.
\end{tcolorbox}
\section[\texorpdfstring{\bn{কড়া} (kôṛā)}{কড়া (kôṛā)}]{\bn{কড়া} (kôṛā) — strict; strong}

\label{lesson:BN-C161-kora}



\begin{quote}
Before the new one: say the Bengali for polite, then the Bengali for gentle.
\end{quote}

\subsection*{You'll want to know: \bn{কড়া}}
\textbf{\bn{কড়া}} — \emph{kôṛā} — \textquotedblleft{}strict; strong\textquotedblright{}.

\begin{grammarlens}[title={What you've built}]
One more word for this chapter.
\end{grammarlens}

\subsection*{Your turn}
\begin{itemize}
\raggedright
  \item \emph{Say it:} \emph{kôṛā}
  \item \emph{Say it:} \emph{kôṛā}, once more
  \item \emph{Say it:} say \emph{kôṛā}
  \item \emph{Recall:} say the Bengali for polite, then the Bengali for gentle, then say \emph{kôṛā} again
\end{itemize}

\begin{tcolorbox}[breakable,colback=teal!4,colframe=teal!35!black,title={Before you move on}]
What does \bn{কড়া} mean? (\textquotedblleft{}Strict; strong\textquotedblright{}.) Say it once more.
\end{tcolorbox}
\section[\texorpdfstring{\bn{সুখী} (sukhī)}{সুখী (sukhī)}]{\bn{সুখী} (sukhī) — contented, happy}

\label{lesson:BN-C161-sukhi}



\begin{quote}
Before the new one: say the Bengali for gentle, then the Bengali for strict; strong.
\end{quote}

\subsection*{You'll want to know: \bn{সুখী}}
\textbf{\bn{সুখী}} — \emph{sukhī} — \textquotedblleft{}contented, happy\textquotedblright{}.

\begin{grammarlens}[title={What you've built}]
One more word for this chapter.
\end{grammarlens}

\subsection*{Your turn}
\begin{itemize}
\raggedright
  \item \emph{Say it:} \emph{sukhī}
  \item \emph{Say it:} \emph{sukhī}, once more
  \item \emph{Say it:} say \emph{sukhī}
  \item \emph{Recall:} say the Bengali for gentle, then the Bengali for strict; strong, then say \emph{sukhī} again
\end{itemize}

\begin{tcolorbox}[breakable,colback=teal!4,colframe=teal!35!black,title={Before you move on}]
What does \bn{সুখী} mean? (\textquotedblleft{}Contented, happy\textquotedblright{}.) Say it once more.
\end{tcolorbox}
\section[\texorpdfstring{\bn{বিরক্ত} (birôktô)}{বিরক্ত (birôktô)}]{\bn{বিরক্ত} (birôktô) — annoyed}

\label{lesson:BN-C161-birokto}



\begin{quote}
Before the new one: say the Bengali for strict; strong, then the Bengali for contented.
\end{quote}

\subsection*{You'll want to know: \bn{বিরক্ত}}
\textbf{\bn{বিরক্ত}} — \emph{birôktô} — \textquotedblleft{}annoyed\textquotedblright{}.

\begin{grammarlens}[title={What you've built}]
One more word for this chapter.
\end{grammarlens}

\subsection*{Your turn}
\begin{itemize}
\raggedright
  \item \emph{Say it:} \emph{birôktô}
  \item \emph{Say it:} \emph{birôktô}, once more
  \item \emph{Say it:} say \emph{birôktô}
  \item \emph{Recall:} say the Bengali for strict; strong, then the Bengali for contented, then say \emph{birôktô} again
\end{itemize}

\begin{tcolorbox}[breakable,colback=teal!4,colframe=teal!35!black,title={Before you move on}]
What does \bn{বিরক্ত} mean? (\textquotedblleft{}Annoyed\textquotedblright{}.) Say it once more.
\end{tcolorbox}
